% GENERATED FILE. Edit canonical lessons, then run npm run generate:books.
% canonical-source-hash: fnv1a64:1e7eecca009e5884
% canonical-lessons: TE-C221-godhuma-rangu, TE-C221-vendi, TE-C221-leta, TE-C221-gundrani, TE-C221-padunaina

\chapter{Brown, Silver, Pale, Round, Sharp}
\label{ch:te-brown-silver-pale-round-sharp}
\hlchaptermodality{221}

\begin{chapteropening}
\textbf{By the end of this chapter:} \emph{I can say brown, silver, pale, round and sharp.}

Complete the last lesson of chapter 221: I can say brown, silver, pale, round and sharp.
\end{chapteropening}

\section[\texorpdfstring{\te{గోధుమ} \te{రంగు} (gōdhuma raṅgu)}{గోధుమ రంగు (gōdhuma raṅgu)}]{\te{గోధుమ} \te{రంగు} (gōdhuma raṅgu) — brown}

\label{lesson:TE-C221-godhuma-rangu}



\begin{quote}
Before the new one: say the Telugu for thin, then the Telugu for raw.
\end{quote}

\subsection*{You'll want to know: \te{గోధుమ} \te{రంగు}}
\textbf{\te{గోధుమ} \te{రంగు}} — \emph{gōdhuma raṅgu} — \textquotedblleft{}brown\textquotedblright{}.

\begin{grammarlens}[title={What you've built}]
One more word for this chapter.
\end{grammarlens}

\subsection*{Your turn}
\begin{itemize}
\raggedright
  \item \emph{Say it:} \emph{gōdhuma raṅgu}
  \item \emph{Say it:} \emph{gōdhuma raṅgu}, once more
  \item \emph{Say it:} say \emph{gōdhuma raṅgu}
  \item \emph{Recall:} say the Telugu for thin, then the Telugu for raw, then say \emph{gōdhuma raṅgu} again
\end{itemize}

\begin{tcolorbox}[breakable,colback=teal!4,colframe=teal!35!black,title={Before you move on}]
What does \te{గోధుమ} \te{రంగు} mean? (\textquotedblleft{}Brown\textquotedblright{}.) Say it once more.
\end{tcolorbox}
\section[\texorpdfstring{\te{వెండి} (veṇḍi)}{వెండి (veṇḍi)}]{\te{వెండి} (veṇḍi) — silver}

\label{lesson:TE-C221-vendi}



\begin{quote}
Before the new one: say the Telugu for raw, then the Telugu for brown.
\end{quote}

\subsection*{You'll want to know: \te{వెండి}}
\textbf{\te{వెండి}} — \emph{veṇḍi} — \textquotedblleft{}silver\textquotedblright{}.

\begin{grammarlens}[title={What you've built}]
One more word for this chapter.
\end{grammarlens}

\subsection*{Your turn}
\begin{itemize}
\raggedright
  \item \emph{Say it:} \emph{veṇḍi}
  \item \emph{Say it:} \emph{veṇḍi}, once more
  \item \emph{Say it:} say \emph{veṇḍi}
  \item \emph{Recall:} say the Telugu for raw, then the Telugu for brown, then say \emph{veṇḍi} again
\end{itemize}

\begin{tcolorbox}[breakable,colback=teal!4,colframe=teal!35!black,title={Before you move on}]
What does \te{వెండి} mean? (\textquotedblleft{}Silver\textquotedblright{}.) Say it once more.
\end{tcolorbox}
\section[\texorpdfstring{\te{లేత} (lēta)}{లేత (lēta)}]{\te{లేత} (lēta) — pale (of a colour); tender, young}

\label{lesson:TE-C221-leta}



\begin{quote}
Before the new one: say the Telugu for brown, then the Telugu for silver.
\end{quote}

\subsection*{You'll want to know: \te{లేత}}
\textbf{\te{లేత}} — \emph{lēta} — \textquotedblleft{}pale (of a colour); tender, young\textquotedblright{}.

\begin{grammarlens}[title={What you've built}]
One more word for this chapter.
\end{grammarlens}

\subsection*{Your turn}
\begin{itemize}
\raggedright
  \item \emph{Say it:} \emph{lēta}
  \item \emph{Say it:} \emph{lēta}, once more
  \item \emph{Say it:} say \emph{lēta}
  \item \emph{Recall:} say the Telugu for brown, then the Telugu for silver, then say \emph{lēta} again
\end{itemize}

\begin{tcolorbox}[breakable,colback=teal!4,colframe=teal!35!black,title={Before you move on}]
What does \te{లేత} mean? (\textquotedblleft{}Pale (of a colour); tender, young\textquotedblright{}.) Say it once more.
\end{tcolorbox}
\section[\texorpdfstring{\te{గుండ్రని} (guṇḍrani)}{గుండ్రని (guṇḍrani)}]{\te{గుండ్రని} (guṇḍrani) — round}

\label{lesson:TE-C221-gundrani}



\begin{quote}
Before the new one: say the Telugu for silver, then the Telugu for pale.
\end{quote}

\subsection*{You'll want to know: \te{గుండ్రని}}
\textbf{\te{గుండ్రని}} — \emph{guṇḍrani} — \textquotedblleft{}round\textquotedblright{}.

\begin{grammarlens}[title={What you've built}]
One more word for this chapter.
\end{grammarlens}

\subsection*{Your turn}
\begin{itemize}
\raggedright
  \item \emph{Say it:} \emph{guṇḍrani}
  \item \emph{Say it:} \emph{guṇḍrani}, once more
  \item \emph{Say it:} say \emph{guṇḍrani}
  \item \emph{Recall:} say the Telugu for silver, then the Telugu for pale, then say \emph{guṇḍrani} again
\end{itemize}

\begin{tcolorbox}[breakable,colback=teal!4,colframe=teal!35!black,title={Before you move on}]
What does \te{గుండ్రని} mean? (\textquotedblleft{}Round\textquotedblright{}.) Say it once more.
\end{tcolorbox}
\section[\texorpdfstring{\te{పదునైన} (padunaina)}{పదునైన (padunaina)}]{\te{పదునైన} (padunaina) — sharp}

\label{lesson:TE-C221-padunaina}



\begin{quote}
Before the new one: say the Telugu for pale, then the Telugu for round.
\end{quote}

\subsection*{You'll want to know: \te{పదునైన}}
\textbf{\te{పదునైన}} — \emph{padunaina} — \textquotedblleft{}sharp\textquotedblright{}.

\begin{grammarlens}[title={What you've built}]
One more word for this chapter.
\end{grammarlens}

\subsection*{Your turn}
\begin{itemize}
\raggedright
  \item \emph{Say it:} \emph{padunaina}
  \item \emph{Say it:} \emph{padunaina}, once more
  \item \emph{Say it:} say \emph{padunaina}
  \item \emph{Recall:} say the Telugu for pale, then the Telugu for round, then say \emph{padunaina} again
\end{itemize}

\begin{tcolorbox}[breakable,colback=teal!4,colframe=teal!35!black,title={Before you move on}]
What does \te{పదునైన} mean? (\textquotedblleft{}Sharp\textquotedblright{}.) Say it once more.
\end{tcolorbox}
